\section{Línea principal de los experimentos: del puntaje global a mecanismos falsables}

La parte experimental no es una tabla de posiciones de «DT ganó», sino una serie de pruebas que se van haciendo más estrictas: primero se comparan los resultados globales de offline RL; después se prueba si la target return es controlable, si el método solo selecciona trayectorias de puntaje alto, si la history es necesaria, si hay degradación con sparse reward, cómo se comporta el credit de largo plazo y si el mismo modelo puede funcionar como critic. Cada figura responde una pregunta distinta.

\subsection{Alcance experimental y protocolo de evaluación}

El trabajo original usa entornos como Atari, OpenAI Gym / MuJoCo y Key-to-Door, y compara con baselines como model-free offline RL y behavior cloning. Cada benchmark define de manera distinta la cantidad de datos, la reward scale y el normalized score; por ello, la comparación debe hacerse dentro del mismo conjunto de datos y del mismo protocolo, y no se pueden tomar directamente los valores de tablas distintas como un porcentaje unificado.

\lecturefigure{slide-14-experiments-divider.jpg}{La clase pasa del mecanismo del método a los experiments: cada figura siguiente pone a prueba una hipótesis concreta.}{Video oficial de Stanford Online, 32:24.}

\begin{knowledgebox}{La intuición del normalized score}
D4RL suele transformar la raw return en
\[
\mathrm{score}_{\mathrm{norm}}
=100\times\frac{R_{\mathrm{policy}}-R_{\mathrm{random}}}
{R_{\mathrm{expert}}-R_{\mathrm{random}}}.
\]
$R_{\mathrm{policy}}$ es el retorno real en el entorno de la policy evaluada; $R_{\mathrm{random}}$ y $R_{\mathrm{expert}}$ son puntos de referencia. Este valor facilita la comparación dentro de una misma tarea, pero no necesariamente cae de forma estricta en $[0,100]$ ni es un «porcentaje de tarea completada».
\end{knowledgebox}

\subsection{Resultados globales de offline RL}

El primer grupo de figuras compara Decision Transformer, los métodos de tipo TD-learning y behavior cloning. Al leer la figura conviene revisar primero tres cosas: a qué environment/data regime corresponde cada grupo de barras, si los baselines usan el mismo offline dataset y si las diferencias son consistentes entre tareas. La conclusión de la clase es que DT es competitivo en varios benchmarks, no que domine de forma absoluta en todas las configuraciones.

\lecturefigure{slide-15-offline-rl-summary-results.jpg}{Offline RL summary: comparación de Decision Transformer con TD learning y behavior cloning en varios grupos de tareas.}{Video oficial de Stanford Online, 33:24; los valores concretos son los de las tablas del artículo original.}

\begin{knowledgebox}{Lectura de la figura: la gráfica de barras respalda «competitive», no «universally superior»}
La altura relativa de los tres grupos de tareas indica que DT puede alcanzar el nivel de baselines fuertes sin entrenar explícitamente una value function. El número de tareas de la figura es limitado, los datos provienen de un collection process predeterminado y los distintos métodos pueden ser sensibles a los hyperparameters y a la architecture. Por ello, la conclusión sólida es «sequence modeling es una offline-RL formulation viable», y no «la programación dinámica ya no sirve».
\end{knowledgebox}

\teachervoice{En la discusión de clase se preguntó si los baselines incluyen recurrence y con qué policy se recolectaron los datos. El expositor reconoce que algunos detalles de implementación deben verificarse y subraya que el replay buffer suele provenir del proceso de entrenamiento de un agent en línea. Experiencia práctica: cuando los resultados dependen del data regime, la source policy y la cantidad de datos deben reportarse junto con la estructura del modelo.}

\subsection{¿El return conditioning es realmente controlable?}

Si el modelo recibe como entrada una desired return, la verificación más directa es fijar distintos targets en el eje horizontal y medir en el eje vertical la achieved return real. El oracle ideal es la diagonal $y=x$; la mejor trayectoria del logged dataset ofrece otra cota superior de referencia. Cuanto más se acerque la curva a la diagonal, más consistentes son la condición de objetivo y la calidad real del comportamiento.

\lecturefigure{slide-16-evaluating-return-conditioning.jpg}{Return conditioning: el eje horizontal es la requested return; el eje vertical, la return obtenida realmente en el rollout.}{Video oficial de Stanford Online, 38:12.}

\begin{knowledgebox}{Lectura de la figura: primero hay que distinguir las tres «cotas superiores»}
La diagonal verde representa el control perfecto: se obtiene exactamente lo que se pide; la línea naranja representa la mejor trayectoria de los datos fuera de línea, no el óptimo teórico del entorno; la curva del modelo representa el rollout real. En algunos entornos, el modelo muestra indicios de mejora por encima del retorno máximo de los datos de entrenamiento, pero con targets más altos termina por saturarse. La saturación es más segura que una caída catastrófica, pero aun así muestra que el target token no es una perilla de ganancia ilimitada.
\end{knowledgebox}

\begin{warningbox}{La occasional extrapolation no es una garantía OOD}
En entornos como Seaquest se ha observado que la achieved return supera la logged best trajectory; esto puede deberse a la combinación de comportamientos locales ya existentes, a la aleatoriedad o a la generalización de la función. El expositor afirma explícitamente que esta tendencia no es consistente. En ingeniería deben reportarse el target sweep, los intervalos de confianza, la tasa de fallos y la distancia respecto del soporte de los datos; no basta con mostrar el mejor intento.
\end{warningbox}

\teachervoice{Un estudiante pregunta: ¿qué pasa si se introduce un target absurdamente grande (por ejemplo, 10,000)? El expositor responde que el desempeño suele saturarse a partir de cierto umbral, en lugar de mejorar sin límite, y que tampoco se ha observado de manera general un colapso repentino. Esta explicación oral es más precisa que las palabras «generación controlable».}

\subsection{«Multitask» requiere acotar su significado}

Cambiar la target return permite que un mismo modelo genere comportamientos distintos, como «quedarse quieto de forma conservadora», «moverse de forma moderada» o «buscar un puntaje alto», por lo que puede considerarse una forma de behavior selection. Pero esto no equivale automáticamente a generalización entre tareas: la return es un escalar unidimensional y puede no distinguir objetivos, restricciones o tareas semánticas diferentes. Un multitask agent más general necesita además un goal state, una language instruction, un task ID u otras condiciones.

\begin{warningbox}{El return conditioning no es una task specification completa}
Dos tareas pueden tener la misma return scale y exigir comportamientos totalmente distintos; una misma return también puede corresponder a varias estrategias. Usar la return como task cue funciona en entornos específicos, pero no prueba por sí solo que el modelo haya aprendido una identidad abstracta de tarea.
\end{warningbox}

\subsection{¿Es solo behavior cloning de trayectorias de puntaje alto?}

Percent behavior cloning (percent BC) primero ordena las trayectorias por trajectory return, conserva el $q\%$ superior y luego entrena un behavior cloning ordinario:
\[
\mathcal{D}_{q}=\{\tau\in\mathcal{D}:R(\tau)\ge Q_{1-q}(R)\},
\qquad
\min_{\theta}\;\mathbb{E}_{(s,a)\sim\mathcal{D}_q}[-\log\pi_{\theta}(a\mid s)].
\]
$Q_{1-q}(R)$ es el umbral de cuantil de la return. Este baseline es fuerte porque también retira de la señal de supervisión las trayectorias de baja calidad; la diferencia es que DT conserva todos los datos y usa la condición de return para distinguir la calidad.

\lecturefigure{slide-17-comparison-with-behavior-cloning.jpg}{Comparación con percent behavior cloning: es muy fuerte en los regímenes de datos medio y alto; en el regime de pocos datos, la ventaja de DT es más evidente.}{Video oficial de Stanford Online, 47:34.}

\begin{knowledgebox}{Lectura de la figura: por qué importa un baseline sencillo}
La tabla compara primero DT con BC usando distintas proporciones de filtrado. Con cantidades de datos medias y altas, percent BC puede acercarse a DT, lo que indica que parte de la ganancia sí proviene de «concentrarse en el comportamiento de alto retorno»; en el regime de pocos datos de Atari, filtrar en exceso deja todavía menos datos, y la ventaja de DT, que aprovecha todas las trayectorias y la condición de return, es más evidente. Un baseline sencillo y fuerte evita atribuir erróneamente a una arquitectura compleja un efecto de selección de datos.
\end{knowledgebox}

\teachervoice{El expositor afirma que percent BC es un baseline fuerte que todo trabajo futuro de offline RL debería incluir. El docente enfatiza que esto no debilita a Decision Transformer, sino que precisa mejor su aportación: cuando hay datos suficientes, filtrar trayectorias de alta calidad ya es muy fuerte; cuando los datos escasean, condicionar sobre todos los datos tiene más valor.}

\subsection{¿La context length es realmente útil?}

Si se fija la context length en $K=1$, el modelo depende casi exclusivamente de la return y el state actuales; si el context completo es claramente mejor, significa que la historia sí aporta action-relevant information. Este ablation pone a prueba directamente el diseño «intencionalmente non-Markov» descrito antes, en lugar de limitarse a comparar redes más grandes y más pequeñas.

\lecturefigure{slide-18-effect-of-context-length.jpg}{Context-length ablation: la historia completa obtiene mejores resultados que $K=1$ en varias tareas.}{Video oficial de Stanford Online, 53:06.}

\begin{knowledgebox}{Lectura de la figura: la diferencia de desempeño puede tener tres orígenes}
La ventana de historia puede recuperar una partial observation, identificar la behavior phase o, simplemente, aportar el episode time. Para confirmar el mecanismo real habría que comparar además: una time feature explícita, un recurrent baseline, una historia barajada al azar, curvas con distintos $K$ y la generalización entre episodios. La slide original respalda que «la historia es útil», pero todavía no permite determinar de manera única qué tipo de información de la historia es la más importante.
\end{knowledgebox}

\subsection{Sparse reward: credit assignment no es lo mismo que exploration}

El experimento de recompensa dispersa solo entrega la cumulative reward al final de la trajectory. Durante el entrenamiento, la return-to-go de cada posición puede recalcularse a partir del registro completo, de modo que los tokens tempranos siguen recibiendo la información condicional de «si esta trayectoria terminó en éxito o en fracaso». Los métodos tradicionales per-transition, si no tienen un backup o una representation adecuados, pueden tener más dificultad para propagar resultados lejanos.

\lecturefigure{slide-19-sparse-reward-environments.jpg}{Experimento de sparse reward: la reward acumulada solo se entrega al final de la trayectoria; DT mantiene su desempeño, mientras que CQL cae notablemente.}{Video oficial de Stanford Online, 55:12.}

\begin{warningbox}{Este experimento no demuestra que la exploración en línea esté resuelta}
El offline dataset ya contiene trajectories exitosas y fallidas, y el modelo puede ver el resultado completo durante el entrenamiento; en cambio, un cold start en línea puede no tener ni una sola trayectoria exitosa. La return-to-go ayuda a propagar el resultado futuro como condición hacia las acciones tempranas, pero no puede descubrir automáticamente la recompensa si no hay datos de éxito. Credit assignment y exploration son problemas relacionados, pero distintos.
\end{warningbox}

\teachervoice{La discusión de clase ofreció una buena intuición: en el offline training, cada paso se reconecta con la trayectoria correcta mediante la action real del registro, de modo que el modelo no avanza «a ciegas» tras unas acciones aleatorias tempranas, como le ocurre a un agent en línea. El expositor coincide en que, al pasar realmente a un entorno en línea, los detalles de exploración y exploitation deben rediseñarse.}

\subsection{Key-to-Door: credit de largo plazo entre fases}

La tabla de sparse reward solo muestra que, al eliminar las rewards intermedias, el desempeño no colapsa de forma evidente; aun así, podrían estar actuando episodios cortos o pistas locales. Key-to-Door separa deliberadamente la acción temprana que determina el éxito de la reward final: el agent debe primero tomar la llave, luego atravesar una habitación vacía sin información de la tarea y solo al final ve la puerta; únicamente obtiene la binary reward si tiene la llave y logra abrir la puerta. Este diseño de tres fases separa las acciones causalmente relevantes de un largo distractor interval, y pone a prueba específicamente si el modelo puede relacionar el éxito final con la toma de la llave ocurrida mucho antes.

\lecturefigure{slide-20-long-term-credit-assignment.jpg}{Key-to-Door: las tres fases (tomar la llave, la habitación vacía y abrir la puerta) ponen a prueba el credit assignment de largo plazo.}{Video oficial de Stanford Online, 1:00:04.}

\begin{knowledgebox}{Lectura de la figura: el trajectory count también es una variable independiente}
La tabla compara distintas cantidades de trajectories exitosas. DT y percent BC superan a varios baselines en esta tarea, pero cuando el registro contiene muy pocos comportamientos exitosos, cualquier método carece de evidencia. El resultado respalda que «el sequence context puede conservar información entre fases», pero no demuestra que el modelo pueda inventar la estrategia de tomar la llave sin ningún ejemplo de éxito.
\end{knowledgebox}

\subsection{¿Puede el mismo modelo convertirse en critic?}

Si se cambia el target de salida de la action a la return o a la reward probability, el causal Transformer puede interpretarse como un critic. Tras observar la historia, el modelo estima la probabilidad de éxito futuro: si no ha tomado la llave, la probabilidad baja; después de tomarla, se mantiene por un tiempo; al acercarse finalmente a la puerta, sube o baja según la acción.

\lecturefigure{slide-21-decision-transformers-as-critics.jpg}{Decision Transformer como critic: actualiza la future reward probability a lo largo de las tres fases de Key-to-Door.}{Video oficial de Stanford Online, 1:03:02.}

\begin{knowledgebox}{Lectura de la figura: los cambios de la curva corresponden a información reversible e irreversible}
En cuanto termina la phase 1 sin que el agent haya tomado la llave, el éxito futuro es prácticamente inalcanzable, y la línea azul debe bajar y mantenerse baja; tras tomar la llave, la habitación vacía de la phase 2 no decide de inmediato el éxito o el fracaso, por lo que las trayectorias candidatas siguen cercanas entre sí; la bifurcación definitiva ocurre en la phase 3, según se llegue o no a la puerta. Esto indica que el hidden state del modelo conserva eventos de largo plazo y puede usarlos para una value-like prediction.
\end{knowledgebox}

\begin{warningbox}{Un «impressive critic» todavía no es un algoritmo actor-critic completo}
Este experimento muestra que el target es intercambiable y que hay capacidad de predicción de largo plazo, pero no aporta automáticamente policy improvement, pessimism, uncertainty calibration ni off-policy correction. Para integrarlo en un sistema actor-critic, aún hay que definir el target, el bootstrap, el data support y la estrategia de estabilización.
\end{warningbox}

\subsection{Resumen de la sección}

La cadena experimental demuestra paso a paso que: DT es una offline-RL formulation competitiva; la return condition y la achieved return guardan una controlabilidad relativa; percent BC es un baseline fuerte con el que es obligatorio comparar; la history sí es útil en algunas tareas; el sequence context puede manejar credit disperso y de largo plazo, y la misma arquitectura puede convertirse en critic. Todas las conclusiones dependen del data regime, y ninguna resuelve directamente la online exploration.
